\section{Analisis Permasalahan dan Asumsi Desain}
\subsection{Lingkup sistem, pengguna, dan masukan}
SmartTrafficLightSystem mengatur satu simpang dengan empat pendekat, yaitu Utara, Selatan, Timur, dan Barat. Model hanya mencakup gerak lurus. Fase pejalan kaki, fase belok khusus, prioritas darurat, dasbor, pembelajaran mesin, dan koordinasi antarsimpang berada di luar lingkup. Setiap pasangan pendekat dianggap berkonflik, sehingga hanya satu pendekat boleh mendapat hijau pada satu waktu. TrafficDataSource menyediakan hitungan kendaraan dan RoadUsers menerima sinyal lampu.

VehicleDetection menerima dan memeriksa data hitungan. Setiap \texttt{CountsFrame} memuat empat catatan dengan ID pendekat yang berbeda, hitungan bilangan bulat dari 0 sampai 2.147.483.647, waktu observasi pada jam monotonik yang sama, dan penanda validitas. Hitungan menunjukkan jumlah kendaraan dalam antrean saat jendela observasi berakhir. Rancangan memakai hitungan itu sebagai indikator kepadatan karena data kapasitas dan luas pendekat belum tersedia.

\begin{table}[htbp]\centering
\caption{Parameter model berdasarkan asumsi desain ASM-002--008}
\begin{tabular}{|p{.43\linewidth}|p{.44\linewidth}|}
\hline
\textbf{Parameter} & \textbf{Nilai atau aturan} \\ \hline
Kategori LOW / MEDIUM / HIGH & 0--5 / 6--15 / $\geq 16$ kendaraan \\ \hline
Hijau LOW / MEDIUM / HIGH & 10 / 20 / 30 s \\ \hline
Hijau minimum / maksimum & 10 / 30 s \\ \hline
Kuning / semua merah & 3 / 1 s \\ \hline
Interval observasi / keputusan & 1 / 1 s \\ \hline
Batas kedaluwarsa & umur data $>3$ s \\ \hline
Waktu tunggu maksimum & 210 s pada operasi nominal \\ \hline
Fallback / pemulihan & Hijau 10 s menurut giliran / 3 data baru valid dan awal putaran berikutnya \\ \hline
\end{tabular}
\end{table}

Tabel ini menetapkan nilai untuk analisis model. Sebagai pembanding, \textcite{mondal2022} menggunakan siklus 120 s dan memberi contoh durasi kuning 2,5 s. Durasi pada simpang nyata perlu ditentukan berdasarkan geometri dan kondisi setempat. Perubahan parameter pada model memerlukan perhitungan ulang batas waktu tunggu.

\subsection{Prioritas dan batas tunggu}
Setiap pendekat mendapat satu giliran per putaran, termasuk saat hitungannya nol. Pada akhir fase semua merah, pengendali hanya mempertimbangkan pendekat yang belum mendapat giliran. Pengendali memilih kategori kepadatan tertinggi terlebih dahulu. Jika kategorinya sama, pengendali mendahulukan pendekat yang menunggu lebih lama. Jika masih seri, urutan N--S--E--W digunakan mulai dari penunjuk giliran, yang bermula di N. Setelah suatu pendekat menerima hijau, penunjuk maju ke pendekat berikutnya. Durasi hijau yang sedang berjalan tetap berjalan hingga habis meski hitungan berubah.

Waktu tunggu dihitung sejak hijau terakhir untuk pendekat tersebut berakhir, atau sejak penyalaan awal jika belum pernah mendapat hijau. Nilainya nol selama hijau. Sebelum memberi giliran berikutnya, pengendali memeriksa apakah semua pendekat yang masih menunggu dapat dilayani dalam 210 s. Jika urutan nominal tidak memenuhi tenggat, pengendali mencoba \textit{earliest-deadline-first} (EDF). Jika EDF juga tidak layak, sistem meminta SafetyHold.

Waktu tunggu terpanjang terjadi ketika sebuah pendekat dilayani pertama pada satu putaran dan terakhir pada putaran berikutnya. Di antara kedua giliran itu ada paling banyak enam fase hijau pendekat lain dan tujuh jeda yang masing-masing mencakup kuning dan semua merah.
\begin{equation}\label{eq:bound}
W_{\max}\leq 6G_{\max}+7(Y+A)=6(30)+7(3+1)=208\ \mathrm{s}<210\ \mathrm{s}.
\end{equation}
Saat penyalaan awal, pendekat terakhir dapat menunggu hingga $1+3(30+4)=103$ s. Selisih 2 s antara batas 208 s dan tenggat 210 s berlaku pada model dengan waktu diskrit. Latensi perangkat belum dimasukkan dalam perhitungan.

\subsection{Penanganan gangguan data dan batas jaminan}
Jika satu sampel hilang, adapter dapat menggunakan data tersimpan yang masih lengkap dan berumur paling lama 3 s. Adapter memeriksanya setiap detik, termasuk ketika masukan berhenti. Data tidak valid, tidak lengkap, atau kedaluwarsa mengaktifkan \texttt{dataFault}. Hitungan yang salah tidak diubah menjadi nol.

Gangguan data tidak memotong fase hijau yang sedang berjalan. Setelah hijau, kuning, dan semua merah selesai, mode Fallback menggunakan urutan N--S--E--W untuk melayani pendekat yang masih menunggu dengan hijau 10 s. Posisi dalam putaran dan waktu tunggu tetap tercatat. Penanda gangguan hilang setelah tiga data lengkap dengan waktu observasi baru diterima berturut-turut pada detik yang berbeda. Mode normal baru dapat dipakai pada awal putaran berikutnya.

SignalController menolak perintah di luar domain atau yang melewati jeda pengaman. Dalam kondisi itu, pengendali masuk SafetyHold dan mempertahankan semua lampu merah hingga dimulai ulang. Batas minimum hijau dan batas waktu tunggu nominal dikecualikan saat SafetyHold atau penyalaan ulang. Analisis ini mengasumsikan jam dan aktuator ideal. Lampu macet, pemadaman, dan latensi belum diuji.
